% !TEX root = ../00_Main/Main.tex
\documentclass[../00_Main/Main.tex]{subfiles}
\begin{document}

Esta tesis compara distintos m\'etodos de aprendizaje por refuerzo en un
problema de asignaci\'on de recursos limitados bajo incertidumbre, con
una pandemia simulada como banco de pruebas. En cada paso llega una
ciudad con un nivel de severidad y el agente decide cu\'antos recursos
le asigna de un presupuesto fijo. La severidad que no se atiende crece
en los pasos siguientes, de modo que una asignaci\'on insuficiente al
principio se paga m\'as adelante. El trabajo se implement\'o en mPES,
una extensi\'on del entorno de pandemia de la Universidad de Essex.

Se evaluaron seis modelos individuales, con hiperpar\'ametros ajustados por optimizaci\'on bayesiana. Dos son m\'etodos tabulares (Q-Learning y Double Q-Learning) y cuatro son redes neuronales: una red profunda (DQN), una variante
recurrente que recuerda los \'ultimos estados (DQN recurrente, con una capa LSTM), un m\'etodo
actor--cr\'itico (A2C) y un Transformer causal, que atiende a la historia
reciente de la secuencia. Tambi\'en se evaluaron seis ensambles que
combinan las decisiones de varias de esas redes. Todos se probaron sobre
las mismas $64$ secuencias, primero en las condiciones de entrenamiento y despu\'es en
$21$ escenarios de generalizaci\'on que cambian la distribuci\'on de
severidades, la longitud de las secuencias o ambas. Tres de ellos llevan
la severidad o la longitud por encima de las cotas del entorno.

El desempe\~no se mide como la fracci\'on de la mejora posible que logra
el agente en cada secuencia: $0$ equivale a no asignar ning\'un recurso y
$1$ a la asignaci\'on \'optima, calculada conociendo de antemano la
secuencia completa. Un agente que decide al azar
obtiene $0{,}670$.

El Transformer fue el mejor modelo individual: obtuvo $0{,}927$ en las
condiciones de entrenamiento, frente a $0{,}887$ de Q-Learning, es
decir, redujo en alrededor de un $36\,\%$ la distancia al \'optimo. En los escenarios de generalizaci\'on mantuvo ese nivel
($0{,}930$, frente a $0{,}871$ de Q-Learning). El mejor sistema fue un
ensamble que combina DQN, la red recurrente y el Transformer y agrega dos
reglas fijas (un \emph{prior} de severidad y una cota de seguridad), con
$0{,}939$ en los escenarios de generalizaci\'on. Los peores resultados
fueron los de los m\'etodos tabulares cuando la severidad supera la cota
del entorno ($0{,}76$ y $0{,}79$, frente a $0{,}996$ del Transformer). Para la mayor\'ia de las redes profundas, el caso m\'as dif\'icil fueron las secuencias m\'as largas que las de entrenamiento.

Los ensambles ponderan a cada modelo seg\'un su confianza en cada
decisi\'on, medida como uno menos la entrop\'ia normalizada de su
distribuci\'on de acciones. El \'unico que super\'o al Transformer le da a
\'este la mayor parte del peso, pero en la referencia su decisi\'on final
difiere de la del Transformer en el $61\,\%$ de los casos, principalmente
por el \emph{prior} de severidad; los que s\'olo promedian o votan quedaron
por debajo.

\par\vspace{0.5\baselineskip}
\noindent\textbf{Palabras clave:} aprendizaje por refuerzo, asignaci\'on
secuencial de recursos, Transformer causal, ensambles de modelos,
entrop\'ia de Shannon, generalizaci\'on.

\biblio % Needed for referencing to working when compiling individual subfiles - Do not remove
\end{document}
